\chapter{คลังรหัสและองค์ประกอบสกิลคลาวด์ฟูลสแตกและ Firebase}

ภาคผนวกนี้รวบรวมข้อมูลจำเพาะ โค้ดต้นแบบ และสถาปัตยกรรมระดับองค์กรของสกิล Firebase AI Logic, Data Connect (PostgreSQL), และ Security Rules Auditor ประจำเล่มที่ 6

\section{สกิล: firebase-ai-logic-basics}

\begin{teachingplan}{ข้อมูลจำเพาะและเมทาดาทาของสกิล}
\textbf{ชื่อสกิล:} \texttt{firebase-ai-logic-basics}\\
\textbf{วัตถุประสงค์:} ผสาน Gemini Multimodal AI เข้าสู่เว็บแอปพลิเคชันอย่างปลอดภัย โดยไม่ต้องเปิดเผย API Key\\
\textbf{ฟังก์ชันการทำงาน:} Streaming Responses, JSON Schema Structured Output, Multimodal Ingestion
\end{teachingplan}

\subsection{โค้ดแกนกลาง: การเรียกใช้งาน Multimodal Inference ผ่าน Firebase SDK}

\begin{lstlisting}[language=JavaScript, caption={Client-Side Firebase AI Logic Integration}]
import { getAI, getGenerativeModel } from "firebase/ai";

const ai = getAI(app);
const model = getGenerativeModel(ai, {
  model: "gemini-1.5-flash",
  generationConfig: {
    responseMimeType: "application/json"
  }
});

async function analyzeLabImage(imageFile, contextPrompt) {
  const imagePart = await fileToGenerativePart(imageFile);
  const result = await model.generateContent([
    contextPrompt, 
    imagePart
  ]);
  return JSON.parse(result.response.text());
}
\end{lstlisting}

\section{สกิล: firebase-data-connect}

\begin{teachingplan}{ข้อมูลจำเพาะและเมทาดาทาของสกิล}
\textbf{ชื่อสกิล:} \texttt{firebase-data-connect}\\
\textbf{วัตถุประสงค์:} ออกแบบและเชื่อมต่อฐานข้อมูล Relational PostgreSQL บน Firebase ผ่าน GraphQL\\
\textbf{ฟังก์ชันการทำงาน:} Schema Definition, Table Relations, Auto-generated Type-Safe SDKs
\end{teachingplan}

\subsection{โค้ดแกนกลาง: การประกาศ Schema ด้วย GraphQL}

\begin{lstlisting}[language=JavaScript, caption={Data Connect GraphQL Schema Blueprint}]
# schema.gql
type Experiment @table {
  id: UUID! @default(expr: "uuidV4()")
  title: String!
  description: String
  domain: String!
  createdAt: Timestamp! @default(expr: "requestTime()")
  records: [Record!]! @ref(fields: "id", references: "experimentId")
}

type Record @table {
  id: UUID! @default(expr: "uuidV4()")
  experimentId: UUID!
  measuredValue: Float!
  uncertainty: Float
  timestamp: Timestamp! @default(expr: "requestTime()")
}
\end{lstlisting}

\section{สกิล: firebase-security-rules-auditor}

\begin{teachingplan}{ข้อมูลจำเพาะและเมทาดาทาของสกิล}
\textbf{ชื่อสกิล:} \texttt{firebase-security-rules-auditor}\\
\textbf{วัตถุประสงค์:} ตรวจสอบและประเมินความปลอดภัยของ Firestore Security Rules แบบ Zero-Trust\\
\textbf{ขอบเขตการตรวจสอบ:} Public Write Vulnerability, User Impersonation, Missing Field Validation
\end{teachingplan}

\subsection{โค้ดแกนกลาง: กฎความปลอดภัยระดับองค์กร}

\begin{lstlisting}[language=bash, caption={Production-Grade Firestore Rules}]
rules_version = '2';
service cloud.firestore {
  match /databases/{database}/documents {
    // User Profiles
    match /users/{userId} {
      allow read: if request.auth != null;
      allow write: if request.auth != null && request.auth.uid == userId;
    }
    
    // Physics Lab Data
    match /experiments/{expId} {
      allow read: if true; // Public view
      allow write: if request.auth != null && 
                      request.resource.data.size() < 50000;
    }
  }
}
\end{lstlisting}
